% Einheit 2 von 4 – Addieren und Subtrahieren (bank/rationale-zahlen/e2.jsonl, zusammenbau v0.1)
%% TODO Zweigzeile Teil 1: was hier gelernt wird (Satz in Schülersprache aus dem Einheitstitel) – Einheit 2
\einheitenkopf[e2]{Einheit 2 von 4 · Addieren und Subtrahieren}
\zweigzeile{ab Kl. 6, je nach Buch bis Kl. 7 (am Gymnasium ab Kl. 7) · P10}
\setcounter{aufgabe}{12}

%% TODO Titel: Ich-kann-Satz fehlt in der Bank (Platzhalter: Kettenname) – Nr. 13
%% TODO Anweisung: ein Satz über der ersten Teilaufgabe fehlt in der Bank – Nr. 13
\begin{aufgabe}{Vorzeichen oder Rechenzeichen?}
\begin{teile}
\teil $-4 - 6$ – jedes Minus: Vorzeichen (V) oder Rechenzeichen (R)? \leerfeld
\end{teile}
\end{aufgabe}

%% TODO Titel: Ich-kann-Satz fehlt in der Bank (Platzhalter: Kettenname) – Nr. 14
%% TODO Anweisung: ein Satz über der ersten Teilaufgabe fehlt in der Bank – Nr. 14
\begin{aufgabe}{Zeichen zusammenfassen}
\begin{teile}
\teil $6 - (-5)$ – ohne Klammern? \leerfeld
\end{teile}
\end{aufgabe}

%% TODO Titel: Ich-kann-Satz fehlt in der Bank (Platzhalter: Kettenname) – Nr. 15
%% TODO Anweisung: ein Satz über der ersten Teilaufgabe fehlt in der Bank – Nr. 15
\begin{aufgabe}{Addieren/Subtrahieren}
%% TODO \rechenplatz{n}: Schrittzahl des Grundfalls fehlt in der Bank (nur bei fester Schreibform, 2.3 b)
\begin{teile}
\teil Zeichne den Pfeil: Start bei $-4$, $6$ nach rechts.

\zahlenstrahl[xmin=-8,xmax=8,xstep=1]{}
\teil Berechne: $-5 + 8$ \leerfeld
\teil Berechne: $7 + (-3)$ \leerfeld
\teil Berechne: $6 - (-7)$ \leerfeld
\teil Berechne: $3 - (+8) + (-4)$ \leerfeld
\teil $-6$ und $5$ – wie weit auseinander? \leerfeld
\teil Berechne: $-2{,}5 + 4{,}1$ \leerfeld
\teil Berechne: $-3 + 8 - 15 + 4$ \leerfeld
\teil Berechne die Klammer von $(a + b) : c$ für $a = 3$ und $b = -17$. \leerfeld
\end{teile}
\end{aufgabe}

%% TODO Titel: Ich-kann-Satz fehlt in der Bank (Platzhalter: Kettenname) – Nr. 16
%% TODO Anweisung: ein Satz über der ersten Teilaufgabe fehlt in der Bank – Nr. 16
\begin{aufgabe}{Beträge: gleiche Vorzeichen addieren, verschiedene subtrahieren}
\begin{teile}
\teil $-17 + (-8)$ – gleiche oder verschiedene Vorzeichen? Rechne mit den Beträgen. \leerfeld
\end{teile}
\end{aufgabe}

%% TODO Titel: Ich-kann-Satz fehlt in der Bank (Platzhalter: Kettenname) – Nr. 17
%% TODO Anweisung: ein Satz über der ersten Teilaufgabe fehlt in der Bank – Nr. 17
\begin{aufgabe}{Dezimalzahlen und Brüche}
\begin{teile}
\teil Berechne: $-\frac{3}{4} + \frac{1}{4}$ \leerfeld
\end{teile}
\end{aufgabe}

%% TODO Titel: Ich-kann-Satz fehlt in der Bank (Platzhalter: Kettenname) – Nr. 18
%% TODO Anweisung: ein Satz über der ersten Teilaufgabe fehlt in der Bank – Nr. 18
\begin{aufgabe}{Umkehrung: Rechnung zu Ergebnis finden}
\begin{teile}
\teil $-7 + \square = 2$ – welche Zahl fehlt? \leerfeld
\end{teile}
\end{aufgabe}

%% TODO Titel: Ich-kann-Satz fehlt in der Bank (Platzhalter: Kettenname) – Nr. 19
%% TODO Anweisung: ein Satz über der ersten Teilaufgabe fehlt in der Bank – Nr. 19
\begin{aufgabe}{Addieren/Subtrahieren – Fehler finden, Begründen, Darstellungswechsel, Anwendung}
\begin{teile}
\teil Ole rechnet: $-4 - 7 = -3$. Finde den Fehler.
\teil Setze fort: $4 - 2 = 2$, $4 - 1 = 3$, $4 - 0 = 4$. Begründe damit, warum $4 - (-3)$ dasselbe ist wie $4 + 3$.
\teil Ein Pfeil startet bei $-3$ und endet bei $4$. Welche Rechnung?

\zahlenstrahl[xmin=-5,xmax=5,xstep=1]{-3/Start, 4/Ende}
\teil Ein Aufzug steht in der Tiefgarage auf Ebene $-2$ und fährt $7$ Stockwerke nach oben. Auf welcher Ebene hält er?
\end{teile}
\end{aufgabe}
